% Part: methods
% Chapter: proofs
% Section: inference-patterns

\documentclass[../../../include/open-logic-section]{subfiles}

\begin{document}

\olfileid{mth}{prf}{pat}

\olsection{अनुमानाचे नमुने}

सिद्धता स्वतंत्र अनुमानांनी बनतात. अनुमान करताना आपण सहसा “म्हणून”, “त्यामुळे” किंवा “यावरून” असे शब्द वापरून ते दर्शवतो. त्या अनुमानाचा आधार अनेकदा सिद्धतेत आधीच उपलब्ध असलेल्या एक-दोन गोष्टी असतात: एखादे गृहीतक किंवा आधीच्या अनुमानाने मिळालेला निष्कर्ष. स्पष्टतेसाठी या गोष्टींना नावे किंवा क्रमांक देता येतात आणि अनुमानात कोणती इतर विधाने वापरली आहेत हे दर्शवता येते. नव्या निष्कर्षाची सत्यता आधीच्या गोष्टींवरून \emph{का} मिळते याचे स्पष्टीकरणही अनुमानात अनेकदा असते. सिद्धतांत वारंवार वापरले जाणारे अनुमानाचे काही सामान्य नमुने आहेत; त्यांपैकी काही पुढे पाहू. विगमनाने सिद्धता करण्यासारखे काही नमुने अधिक गुंतागुंतीचे आहेत; त्यांची चर्चा नंतर होईल.

अनुमानाचा एक नमुना आपण आधीच पाहिला आहे: व्याख्या उलगडणे किंवा तिचा वापर करणे. व्याख्या उलगडताना व्याख्येय असलेले म्हणणे आपण व्याख्यापकाच्या रूपात पुन्हा मांडतो. उदाहरणार्थ, सिद्धतेत आधीच $D = E$ हे (a) प्रस्थापित केले आहे असे समजा. मग संचांसाठीच्या $=$ च्या व्याख्येवरून असे अनुमान करता येते: “म्हणून (a) आणि व्याख्येवरून,~$D$ मधील प्रत्येक !!{element} हा~$E$ मधील !!a{element} असतो आणि उलटही तसेच असते.”

थोडा गोंधळ होऊ शकतो: अनेकदा अनुमानाचे समर्थन प्रत्यक्ष अनुमान करताना लिहिण्याऐवजी त्याआधीच लिहिले जाते. $D = E$ अद्याप सिद्ध केलेले नाही, पण ते सिद्ध करायचे आहे असे समजा. अपेक्षित निष्कर्ष $D = E$ असेल, तर व्याख्या वापरून हे ध्येयही पुन्हा मांडता येते: $D = E$ सिद्ध करण्यासाठी~$D$ मधील प्रत्येक !!{element} हा~$E$ मधील !!a{element} आहे आणि उलटही तसेच आहे, हे सिद्ध करावे लागते. म्हणून सिद्धतेचे रूप असे असेल: (a)~$D$ मधील प्रत्येक !!{element} हा~$E$ मधील !!a{element} आहे हे सिद्ध करा; (b)~$E$ मधील प्रत्येक !!{element} हा~$D$ मधील !!a{element} आहे हे सिद्ध करा; (c) म्हणून (a), (b) आणि $=$ च्या व्याख्येवरून $D = E$. पण सहसा आपण ते अशा रीतीने लिहीत नाही. त्याऐवजी असे लिहू शकतो:
\begin{quote}
$D = E$ दाखवायचे आहे. ~$=$ च्या व्याख्येनुसार त्यासाठी~$D$ मधील प्रत्येक !!{element} हा~$E$ मधील !!a{element} आहे आणि उलटही तसेच आहे, हे दाखवणे पुरेसे आहे.

(a) \dots (~$D$ मधील प्रत्येक !!{element} हा~$E$ मधील !!a{element} आहे याची सिद्धता) \dots

(b) \dots (~$E$ मधील प्रत्येक !!{element} हा~$D$ मधील !!a{element} आहे याची सिद्धता) \dots
\end{quote}

\subsection{संधियोगाचा वापर}

संधियोगातील एखादा संधिघटक निष्कर्ष म्हणून काढणे हा कदाचित अनुमानाचा सर्वात साधा नमुना आहे. म्हणजे $p$ आणि~$q$ असे गृहीत धरले असेल किंवा आधीच सिद्ध केले असेल, तर~$p$ चे (आणि~$q$ चेही) अनुमान करता येते. हे इतके मूलभूत अनुमान आहे की अनेकदा त्याचा स्वतंत्र उल्लेख होत नाही. उदाहरणार्थ, $D = E$ ची व्याख्या उलगडल्यावर~$D$ मधील प्रत्येक !!{element} हा~$E$ मधील !!a{element} आहे आणि उलटही तसेच आहे, हे आपल्याला मिळते. यावरून~$E$ मधील प्रत्येक !!{element} हा~$D$ मधील !!a{element} आहे असा निष्कर्ष काढता येतो; तोच “उलटही तसेच” हा भाग आहे.

\subsection{संधियोग सिद्ध करणे}

कधी सिद्ध करावयाचे विधान संधियोगाच्या रूपात असेल: “$p$ आणि $q$ सिद्ध करा” असे सांगितले जाईल. अशा वेळी दोन गोष्टी कराव्या लागतात: $p$ सिद्ध करा आणि नंतर $q$ सिद्ध करा. सिद्धता दोन भागांत विभागून स्पष्टतेसाठी त्या भागांना नावे किंवा क्रमांक देता येतात. सुरुवातीच्या टिपणांत पानाच्या वरच्या बाजूला “(1) $p$ सिद्ध करा” आणि मध्यभागी “(2) $q$ सिद्ध करा” असे लिहिता येईल. अर्थात, संधियोग सिद्ध करायला स्पष्टपणे सांगितलेले नसले तरी सिद्धतेत संधियोग सिद्ध करणे आवश्यक ठरू शकते. उदाहरणार्थ, $D = E$ सिद्ध करायला सांगितले असेल, तर~$=$ ची व्याख्या उलगडल्यावर~$D$ मधील प्रत्येक !!{element} हा~$E$ मधील !!a{element} आहे \emph{आणि}~$E$ मधील प्रत्येक !!{element} हा~$D$ मधील !!a{element} आहे, हे सिद्ध करावे लागते असे दिसेल.

\subsection{विकल्पयोग सिद्ध करणे}

सिद्ध करावयाचे विधान विकल्पयोगाच्या रूपात असेल---म्हणजे “$p$ किंवा $q$” या रूपाचे असेल---तर एका विकल्पघटकाची सत्यता दाखवणे पुरेसे असते. पण प्रमेयाच्या गृहीतकांवरून कोणताही एक विकल्पघटक आपोआप मिळतो, असे सहसा घडत नाही. उलट, अनेकदा प्रमेयाची गृहीतकेच विकल्पयोगात्मक असतात; किंवा एका प्रकारच्या सर्व वस्तूंत दोनपैकी एक गुणधर्म आहे हे दाखवायचे असते, पण काही वस्तूंत पहिला आणि इतरांत दुसरा गुणधर्म असतो. अशा वेळी प्रकरणांनुसार सिद्धता उपयुक्त ठरते; ती पुढे पाहू.

\subsection{सोपाधिक सिद्धता}

अनेक प्रमेये सशर्त रूपात असतात: $p$ असेल, तर $q$ देखील सत्य आहे, हे दाखवायचे असते. अशा सिद्धतेची मांडणी सहज करता येते: सशर्त विधानाचे पूर्वांग---येथे $p$---गृहीत धरा आणि त्यावरून निष्कर्ष~$q$ सिद्ध करा. म्हणून प्रमेय “$p$ असेल, तर $q$” असे असेल, तर सिद्धतेची सुरुवात “$p$ असे गृहीत धरा” अशी करा आणि शेवटी~$q$ सिद्ध झालेले असले पाहिजे.

सशर्त विधाने वेगवेगळ्या रीतीने मांडता येतात. “$p$ असेल, तर $q$” याऐवजी प्रमेय “$q$ असेल तरच $p$”, “$p$ असेल तर $q$” किंवा “$p$ या अटीवर $q$” असे म्हणू शकते. या सर्वांचा एकच अर्थ आहे: $p$ गृहीत धरून त्यावरून~$q$ सिद्ध करावे लागते. द्विसशर्त विधान---“$p$ सत्य असते तेव्हाच आणि फक्त तेव्हा जेव्हा $q$ सत्य असते” (इंग्रजी संक्षेप “iff”)---हे दोन सशर्त विधानांचे एकत्रीकरण असते, हे आठवा: $p$ असेल तर $q$ आणि $q$ असेल तर~$p$. म्हणून सोपाधिक सिद्धता दोनदा करावी लागते: एकदा पहिल्या सशर्त विधानासाठी आणि एकदा दुसऱ्यासाठी. मात्र कधी इतर “तेव्हाच आणि फक्त तेव्हा” विधानांची साखळी जोडूनही असे विधान सिद्ध करता येते: सुरुवात “$p$” पासून होते आणि शेवट “$q$” येथे. अशा वेळी प्रत्येक पायरी खरोखर “तेव्हाच आणि फक्त तेव्हा” अशीच आहे याची खात्री केली पाहिजे.

\subsection{वैश्विक विधाने}

वैश्विक विधानाचा वापर सोपा आहे: एखादी गोष्ट प्रत्येक वस्तूसाठी खरी असेल, तर ती प्रत्येक विशिष्ट वस्तूसाठी खरी असते. उदाहरणार्थ, सिद्धतेचे गृहीतक $A \subseteq B$ असेल, तर~$\subseteq$ ची व्याख्या उलगडल्यावर त्याचा अर्थ असा होतो: प्रत्येक $x \in A$ साठी $x \in B$ असते. त्यामुळे $z \in A$ हे आधीच माहीत असेल, तर $z \in B$ असा निष्कर्ष काढता येतो.

वैश्विक विधान सिद्ध करणे थोडे अवघड वाटू शकते. सहसा ही विधाने “$x$ मध्ये~$P$ हा गुणधर्म असेल, तर त्यात~$Q$ हाही गुणधर्म असतो” किंवा “सर्व $P$ असलेल्या वस्तू $Q$ असलेल्या असतात” या रूपात असतात. अर्थात, प्रत्येक विधान नेमके या रूपात असेलच असे नाही; नेमके काय सिद्ध करायला सांगितले आहे हे ओळखण्यासाठी थोडा सराव लागतो. पण अनेकदा एखादा गुणधर्म असलेल्या सर्व वस्तूंमध्ये आणखी एक विशिष्ट गुणधर्म आहे, हे सिद्ध करावे लागते.

वैश्विक विधान सिद्ध करण्यासाठी पहिला गुणधर्म असलेल्या वस्तूंना नावे किंवा चल द्यायचे आणि त्यांत दुसराही गुणधर्म आहे हे दाखवायचे. हे असेही म्हणता येईल: \emph{सर्व}~$P$ असलेल्या वस्तूंसाठी काहीतरी सिद्ध करायचे असेल, तर एका \emph{स्वेच्छ}~$P$ असलेल्या वस्तूसाठी ते सिद्ध करावे लागते. दिलेले नाव हे अशा स्वेच्छ~$P$ वस्तूचे नाव असते. स्वेच्छ वस्तूंसाठी सहसा एकेक अक्षर वापरतो आणि त्या अक्षरांसाठी काही रूढी असतात: नैसर्गिक संख्यांसाठी $n$, !!{formula}s साठी $!A$, संचांसाठी $A$, फलनांसाठी $f$, इत्यादी.

$P$ या गुणधर्मापलीकडे विशेष अटी न घालता संपूर्ण सिद्धतेत सामान्यता टिकवणे हा मुख्य मुद्दा आहे. एका स्वेच्छ वस्तूत (“$x$”)~$P$ हा गुणधर्म आहे असे गृहीत धरून सुरुवात करायची आणि केवळ व्याख्या किंवा अनुमत गृहीतकांच्या आधारे $x$ मध्ये~$Q$ हा गुणधर्म आहे हे दाखवायचे. $x$ नेमके काय आहे याविषयी आणखी कोणतीही विशेष अट घातलेली नसल्यामुळे, $P$ असलेल्या प्रत्येक वस्तूत~$Q$ आहे असे म्हणता येते. म्हणजे $x$ हे~$P$ गुणधर्म असलेल्या \emph{सर्व} वस्तूंचे प्रतिनिधी नाव आहे.

\begin{prop}
सर्व $A$ आणि $B$ संचांसाठी, $A \subseteq A \cup B$.
\end{prop}

\begin{proof}
$A$ आणि $B$ हे स्वेच्छ संच असू द्या. $A \subseteq A \cup B$ दाखवायचे आहे. $\subseteq$ च्या व्याख्येनुसार याचा अर्थ असा: प्रत्येक $x$ साठी, $x \in A$ असेल तर $x \in A \cup B$ असते. म्हणून $x \in A$ हा~$A$ मधील स्वेच्छ !!{element} असू द्या. $x \in A \cup B$ दाखवायचे आहे. $x \in A$ असल्यामुळे $x \in A$ किंवा $x \in B$ असते. म्हणून $x \in \Setabs{x}{x \in A \lor x \in B}$. पण $\cup$ च्या व्याख्येनुसार याचाच अर्थ $x \in A \cup B$ असा होतो.
\end{proof}

\subsection{प्रकरणांनुसार सिद्धता}

गृहीतक म्हणून किंवा आधीच मिळालेला निष्कर्ष म्हणून विकल्पयोग उपलब्ध आहे असे समजा: $p$ किंवा $q$ सत्य आहे असे गृहीत धरलेले आहे किंवा सिद्ध झालेले आहे. तुम्हाला $r$ सिद्ध करायचे आहे. हे दोन पायऱ्यांत करता येते: प्रथम $p$ सत्य आहे असे गृहीत धरून~$r$ सिद्ध करा; नंतर $q$ सत्य आहे असे गृहीत धरून पुन्हा~$r$ सिद्ध करा. दोन पर्यायांपैकी एक खरा आहे असे गृहीत धरलेले किंवा माहीत असल्यामुळे ही पद्धत चालते. या दोन पायऱ्यांमुळे कोणताही एक पर्याय~$r$ च्या सत्यतेसाठी पुरेसा आहे हे मिळते. दोन्ही पर्याय खरे असतील, तर $r$~खरे असण्याची एकाऐवजी दोन कारणे असतील. $p$ आणि~$q$ दोन्ही गृहीत धरून $r$~सत्य आहे हे वेगळे सिद्ध करण्याची गरज नाही. काय करत आहोत हे दर्शवण्यासाठी “प्रकरणे वेगळी पाहू” असे जाहीर करतो. उदाहरणार्थ, $x \in B \cup C$ माहीत आहे असे समजा. $B \cup C$ ची व्याख्या $\Setabs{x}{x \in B \text{ किंवा } x \in C}$ अशी आहे. म्हणजे व्याख्येनुसार $x \in B$ किंवा $x \in C$. यावरून $x \in A$ सिद्ध करण्यासाठी प्रथम $x \in B$ गृहीत धरून त्यावरून $x \in A$ सिद्ध करू; नंतर $x \in C$ गृहीत धरून पुन्हा $x \in A$ सिद्ध करू. गृहीतकाखाली “प्रकरणे वेगळी पाहू” लिहा; त्याखाली “प्रकरण (1): $x \in B$” लिहा आणि पानाच्या मध्यावर “प्रकरण (2): $x \in C$” लिहा. मग वरच्या आणि खालच्या भागांतील सिद्धता पूर्ण करा.

सिद्ध करावयाचे विधान स्वतः विकल्पयोगात्मक असेल, तर प्रकरणांनुसार सिद्धता विशेष उपयुक्त ठरते. हे एक साधे उदाहरण:

\begin{prop}
$B \subseteq D$ आणि $C \subseteq E$ असे समजा. मग $B \cup C \subseteq D \cup E$.
\end{prop}

\begin{proof}
(a) $B \subseteq D$ आणि (b) $C \subseteq E$ असे गृहीत धरा. व्याख्येनुसार, कोणताही $x \in B$ हा $\in D$ देखील असतो (c) आणि कोणताही $x \in C$ हा $\in E$ देखील असतो (d). $B \cup C \subseteq D \cup E$ दाखवण्यासाठी $x \in B \cup C$ असेल, तर $x \in D \cup E$ असते हे दाखवावे लागते; हे $\subseteq$ च्या व्याख्येवरून मिळते. ~$ \cup$ च्या व्याख्येनुसार $x \in B \cup C$ हे $x \in B$ किंवा $x \in C$ असते तेव्हाच आणि फक्त तेव्हाच सत्य असते. त्याचप्रमाणे $x \in D \cup E$ हे $x \in D$ किंवा $x \in E$ असते तेव्हाच आणि फक्त तेव्हाच सत्य असते. म्हणून हे दाखवावे लागते: कोणत्याही $x$ साठी, $x \in B$ किंवा $x \in C$ असेल, तर $x \in D$ किंवा $x \in E$ असते.

\begin{quote}
आत्तापर्यंत फक्त व्याख्या उलगडल्या आहेत!{} आपल्या विधानाची $\subseteq$ आणि $\cup$ न वापरता पुनर्मांडणी केली आहे; आता एक वैश्विक सशर्त विधान सिद्ध करायचे उरले आहे. आधीच्या चर्चेनुसार, $x$ ही अशी वस्तू आहे की तिच्याविषयी “जर” भाग सत्य आहे असे गृहीत धरायचे आणि मग “तर” भागही सत्य आहे हे दाखवायचे. म्हणजे $x \in B$ किंवा $x \in C$ असे गृहीत धरून $x \in D$ किंवा $x \in E$ दाखवू.\footnote{हा परिच्छेद आपण काय करत आहोत याचे स्पष्टीकरण देतो; तो सिद्धतेचा भाग नाही. स्वतःच्या सिद्धता लिहिताना हा सगळा तपशील देण्याची गरज नाही.}
\end{quote}

$x \in B$ किंवा $x \in C$ असे समजा. $x \in D$ किंवा $x \in E$ दाखवायचे आहे. प्रकरणे वेगळी पाहू.

प्रकरण 1: $x \in B$. (c) वरून $x \in D$. म्हणून $x \in D$ किंवा $x \in E$. येथे एका विकल्पघटकावरून विकल्पयोग मिळवण्याचे आधी चर्चिलेले अनुमान केले आहे.

प्रकरण 2: $x \in C$. (d) वरून $x \in E$. म्हणून $x \in D$ किंवा $x \in E$.
\end{proof}

\subsection{अस्तित्ववाची विधान सिद्ध करणे}

अस्तित्ववाची विधान सिद्ध करायला सांगितले असेल, तर प्रश्नाचे रूप सहसा “असा~$x$ आहे की $\dots x \dots$, हे सिद्ध करा” असे असते. म्हणजे “$\dots x \dots$” ने वर्णन केलेला गुणधर्म असलेली वस्तू अस्तित्वात आहे, हे दाखवायचे असते. अशा वेळी योग्य वस्तू ओळखून तिच्यात आवश्यक गुणधर्म आहे हे दाखवावे लागते. हे सरळ वाटते; पण अशा सिद्धता अवघडही असू शकतात. सहसा त्यात एखादी वस्तू \emph{रचावी} किंवा \emph{परिभाषित} करावी लागते आणि अशा रीतीने परिभाषित केलेल्या वस्तूत आवश्यक गुणधर्म आहे हे सिद्ध करावे लागते. योग्य वस्तू शोधणे अवघड असू शकते; तिच्यात आवश्यक गुणधर्म आहे हे सिद्ध करणे अवघड असू शकते; आणि कधी आपण मुळात एखादी वस्तू परिभाषित करण्यात यशस्वी झालो आहोत हेच दाखवणे अवघड असते!{}

सहसा वस्तू स्पष्टपणे सांगून अशी सिद्धता लिहितो. उदाहरणार्थ, “$x$ म्हणजे \dots{} असू द्या” असे लिहितो; येथे \dots{} ने कोणती वस्तू अभिप्रेत आहे ते सांगितले जाते. आवश्यक असल्यास $\dots$ खरोखर अस्तित्वात असलेल्या वस्तूचे वर्णन करते हे सिद्ध करतो; मग $x$ मध्ये~$Q$ हा गुणधर्म आहे हे दाखवतो. हे एक साधे उदाहरण:

\begin{prop}
$x \in B$ असे समजा. मग असा~$A$ आहे की $A \subseteq B$ आणि $A \neq \emptyset$.
\end{prop}

\begin{proof}
$x \in B$ असे गृहीत धरा. $A = \{x\}$ असू द्या.
\begin{quote}
येथे~$A$ संच त्याच्या !!{element}s ची यादी देऊन परिभाषित केला आहे. $x$ ही वस्तू आहे असे गृहीत धरले आहे; आणि सांत यादीत !!{element}s देऊन संच तयार करता येतो. म्हणून येथे~$A$ संच परिभाषित करण्यात यश आले आहे हे वेगळे दाखवण्याची गरज नाही. मात्र $A$ मध्ये विधानाने अपेक्षित केलेले गुणधर्म आहेत हे अद्याप दाखवावे लागते. त्याशिवाय सिद्धता पूर्ण होत नाही!{}
\end{quote}
$x \in A$ असल्यामुळे $A \neq \emptyset$.
\begin{quote}
याचा आधार म्हणजे $A$ ची $\{x\}$ अशी व्याख्या आणि $x \in \{x\}$ तसेच $x \notin \emptyset$ या स्पष्ट गोष्टी.
\end{quote}
$x$ हा~$\{x\}$ मधील एकमेव !!{element} आहे आणि $x \in B$ आहे. म्हणून~$A$ मधील प्रत्येक !!{element} हा~$B$ मधील !!a{element} देखील असतो. ~$ \subseteq$ च्या व्याख्येनुसार $A \subseteq B$.
\end{proof}

\subsection{अस्तित्ववाची विधानांचा वापर}

एखादे अस्तित्ववाची विधान सत्य आहे हे माहीत आहे असे समजा: ते सिद्ध केलेले असेल किंवा वापरता येण्याजोगे गृहीतक असेल. उदाहरणार्थ, “काही~$x$ साठी $x \in A$” किंवा “असा $x \in A$ आहे”. सिद्धतेत हे वापरायचे असेल, तर गृहीतकानुसार अस्तित्वात असलेल्या वस्तूंपैकी एका वस्तूला नाव दिले आहे असे धरता येते. $A$ मध्ये किमान एक वस्तू असल्यामुळे त्या नावाने दर्शवता येण्याजोग्या वस्तू आहेत. एखादी वस्तू नेमकी निवडून दाखवता येईलच असे नाही किंवा ती $\in A$ आहे यापलीकडे तिचे वर्णन करता येईलच असे नाही. पण सिद्धतेपुरते एका अशा वस्तूला नाव देता येते. आधी वापरलेले किंवा गृहीतकांत आलेले नाव पुन्हा वापरू नये; तसे केल्यास चुका होऊ शकतात. सिद्धतेत “काही $x$ साठी $x \in A$” यावरून “$a \in A$ असू द्या” असे लिहून हे दर्शवतो. आता~$a$ विषयी विचार करता येतो आणि इतर अनुमत गृहीतके वापरून $p$ हा निष्कर्ष मिळवता येतो. $p$ मध्ये~$a$ चा उल्लेख उरलेला नसेल आणि या नव्या नावाशी संबंधित अतिरिक्त, मुक्त न केलेले गृहीतक वापरलेले नसेल, तर $p$ मिळवण्यासाठीचे तात्पुरते $a \in A$ हे गृहीतक मुक्त करता येते. त्यामुळे “काही $x$ साठी $x \in A$” आणि वापरलेली इतर अनुमत गृहीतके यांवरून निष्कर्ष मिळतो.

\begin{prop}
$A \neq \emptyset$ असेल, तर $A \cup B \neq \emptyset$.
\end{prop}

\begin{proof}
$A \neq \emptyset$ असे समजा. म्हणून काही $x$ साठी $x \in A$.
\begin{quote}
येथे प्रथम विधानाचे गृहीतक फक्त पुन्हा मांडले आहे. या गृहीतकात, म्हणजे $A \neq \emptyset$ यात, एक अस्तित्ववाची विधान दडलेले आहे. काही व्याख्या उलगडल्यावर ते दिसते. $=$ च्या व्याख्येनुसार $A = \emptyset$ हे तेव्हाच आणि फक्त तेव्हा सत्य असते जेव्हा प्रत्येक $x \in A$ हा $\in \emptyset$ देखील असतो आणि प्रत्येक $x \in \emptyset$ हा $\in A$ देखील असतो. दोन्ही बाजूंचे नकरण केल्यावर हे मिळते: $A \neq \emptyset$ तेव्हाच आणि फक्त तेव्हा जेव्हा काही $x \in A$ हा $\notin \emptyset$ असतो किंवा काही $x \in \emptyset$ हा $\notin A$ असतो. कोणतीही वस्तू $\in \emptyset$ नसल्यामुळे दुसरा विकल्पघटक कधीही सत्य असू शकत नाही. तसेच “$x \in A$ आणि $x \notin \emptyset$” याचा अर्थ फक्त $x \in A$ इतकाच उरतो. म्हणून $A \neq \emptyset$ हे काही $x$ साठी $x \in A$ असते तेव्हाच आणि फक्त तेव्हाच सत्य असते. हे अस्तित्ववाची विधान आहे. आता~$A$ मधील !!{element}s पैकी एकाला नाव देऊन या अस्तित्ववाची विधानाचा वापर करू:
\end{quote}
$a \in A$ असू द्या.
\begin{quote}
आता~$\in A$ असलेल्या एका वस्तूला नाव दिले आहे. पुढे~$a$ विषयी युक्तिवाद करू; पण केवळ $a \in A$ आहे एवढेच गृहीत धरू, आणखी काही नाही, याची काळजी घेऊ:
\end{quote}
$a \in A$ असल्यामुळे~$\cup$ च्या व्याख्येवरून $a \in A \cup B$. म्हणून काही $x$ साठी $x \in A \cup B$; म्हणजे $A \cup B \neq \emptyset$.
\begin{quote}
शेवटच्या पायरीत “$a \in A \cup B$” यावरून “काही $x$ साठी $x \in A \cup B$” असा निष्कर्ष काढला. त्यात आता $a$ चा उल्लेख नाही. म्हणून “काही $x$ साठी $x \in A \cup B$” हे केवळ “काही $x$ साठी $x \in A$” यावरून मिळते हे समजते. पण याचाच अर्थ $A \cup B \neq \emptyset$ असा होतो.
\end{quote}
\end{proof}

आपण केले तसे “$x$” सारखे बद्ध चल आणि $a$ सारखी तात्पुरती नावे वेगळी ठेवणे हा चांगला सराव आहे. प्रत्यक्षात मात्र अनेकदा असे करत नाही आणि $x$ हेच वापरतो:
\begin{quote}
$A \neq \emptyset$ असे समजा; म्हणजे असा $x \in A$ आहे. $\cup$ च्या व्याख्येनुसार $x \in A \cup B$. म्हणून $A \cup B \neq \emptyset$.
\end{quote}
मात्र असे करताना वेगवेगळ्या अस्तित्ववाची विधानांसाठी वेगवेगळी $x$ आणि $y$ नावे वापरतो आहोत याची विशेष काळजी घ्यावी. उदाहरणार्थ, पुढील युक्तिवाद ही “$A \neq \emptyset$ आणि $B \neq \emptyset$ असतील, तर $A \cap B \neq \emptyset$” याची योग्य सिद्धता \emph{नाही}; ते विधानच खरे नाही.
\begin{quote}
$A \neq \emptyset$ आणि $B \neq \emptyset$ असे समजा. म्हणून काही $x$ साठी $x \in A$ आणि काही $x$ साठी $x \in B$ देखील असते. $x \in A$ आणि $x \in B$ असल्यामुळे~$\cap$ च्या व्याख्येनुसार $x \in A \cap B$. म्हणून $A \cap B \neq \emptyset$.
\end{quote}
चुकीची पायरी कुठे येते ते ओळखता येईल का? हा निष्कर्ष का खरा नाही हे स्पष्ट करता येईल का?

\end{document}
